\section{Paper Localization Methods}
\label{sec:localization}

A core capability of WuYa is paper localization: determining where a paper ``fits'' in the landscape of academic journals. We propose a dual-path approach where Path A (LLM-as-Mapper) provides fast estimation and Path B (\dea{} Analysis) provides depth, with cross-validation ensuring robustness.

\subsection{Problem Formulation}

Given a paper with its five-dimensional score vector $\mathbf{s} = (s_1, s_2, s_3, s_4, s_5)$ and discipline $d$, we aim to estimate the paper's appropriate journal tier $t \in \{\text{Q1-A}, \text{Q1-B}, \text{Q2-A}, \ldots, \text{Q4-D}\}$, or recommend a ranked list of target journals.

\subsection{Path A: LLM-as-Mapper with Discipline Priors}

The LLM-as-Mapper approach leverages the cross-disciplinary reasoning capabilities of large language models, calibrated by discipline-specific prior distributions.

\subsubsection{Discipline Prior Extraction}

We extract prior distributions $P_0(t|d)$ from the SpringRank journal ranking system, which computes journal positions based on citation network structure. For each discipline $d$, we compute the empirical distribution of journal tiers:
\begin{equation}
P_0(t|d) = \frac{\# \text{journals in tier } t \text{ of discipline } d}{\# \text{total journals in discipline } d}
\end{equation}

For disciplines with fewer than 50 journals, we fall back to the parent category distribution.

\subsubsection{LLM Mapping Prompt}

The LLM receives a carefully designed prompt containing:

\begin{enumerate}
    \item \textbf{Scoring rubric}: The five-dimensional scoring criteria with level descriptions.
    \item \textbf{Few-shot examples}: Representative (score vector, tier) pairs across disciplines.
    \item \textbf{Discipline prior}: The distribution $P_0(t|d)$ formatted as a string.
    \item \textbf{Output template}: Requiring reasoning chain, estimated tier, and confidence level.
\end{enumerate}

\begin{verbatim}
You are a research quality evaluator. Given the following five-dimensional
scores (1=low to 5=high):
- Innovation: [novelty, knowledge_bridge, future_potential, destruction]
- Method: [causal, internal_validity, external_validity, statistical]
- Evidence: [falsifiability, evidence_strength, replicability, program]
- Application: [trl, translation_path, feedback, market]

Discipline: {discipline}
Prior distribution: {P0_string}

Estimate the appropriate journal tier (Q1-A to Q4-D) and explain your reasoning.
\end{verbatim}

\subsubsection{Calibration Strategy}

To prevent LLM hallucination on unfamiliar disciplines, we apply post-hoc calibration:
\begin{equation}
P_{\text{final}}(t|\mathbf{s}, d) \propto \alpha \cdot P_{\text{LLM}}(t|\mathbf{s}, d) + (1-\alpha) \cdot P_0(t|d)
\end{equation}
where $\alpha \in [0.6, 0.8]$ gives primary weight to LLM reasoning while anchoring to factual prior distributions.

To ensure output consistency, we set temperature $\tau = 0.1$, run three generations, and take the modal answer.

\subsection{Path B: DEA Efficiency Analysis}

For deeper analysis, especially when targeting a specific journal, we employ \dea{} efficiency analysis.

\subsubsection{Input-Output Transformation}

We transform the five-dimensional scores into \dea{} inputs (lower is better) and outputs (higher is better):

\begin{itemize}
    \item \textbf{Inputs} (minimize):
    \begin{itemize}
        \item Method flaws: $6 - s_{\text{method}}$
        \item Evidence gap: $6 - s_{\text{evidence}}$
    \end{itemize}
    \item \textbf{Outputs} (maximize):
    \begin{itemize}
        \item Innovation: $s_{\text{innovation}}$
        \item Application: $s_{\text{application}}$
    \end{itemize}
\end{itemize}

This transformation reflects the assumption that strong methods and solid evidence are ``investments'' that should yield innovative and impactful ``returns.''

\subsubsection{Frontier Construction}

For a target journal $j$, we construct the reference frontier from previously published papers in $j$. We require at least 50 papers to build a reliable frontier. If insufficient data exists for $j$, we fall back to a discipline-level frontier constructed from journals in the same tier.

We employ the super-efficiency \dea{} model, which allows efficiency scores greater than 1 for units on the frontier, enabling ranking of even the most efficient papers:
\begin{align}
\max \quad & \theta \\
\text{s.t.} \quad & \sum_{i \in R} \lambda_i x_{ik} \leq \theta x_{k} \\
& \sum_{i \in R} \lambda_i y_{ik} \geq y_{k} \\
& \lambda_i \geq 0
\end{align}
where $R$ is the reference set, $x_k, y_k$ are the input/output of paper $k$, and $\theta > 1$ indicates super-efficiency.

\subsubsection{Bootstrap Confidence Intervals}

To quantify uncertainty, we apply bootstrap re-sampling \citep{efron1994introduction}:

\begin{enumerate}
    \item Resample the reference set with replacement, maintaining original size.
    \item Recompute efficiency scores on the bootstrap frontier.
    \item Repeat 200 times to construct empirical distribution.
    \item Report the 2.5th and 97.5th percentiles as confidence intervals.
\end{enumerate}

A paper whose confidence interval includes 1.0 is considered statistically indistinguishable from the frontier.

\subsection{Cross-Validation and Output}

The outputs of Path A and Path B are compared:

\begin{itemize}
    \item \textbf{Consistent}: Both paths suggest similar tiers $\rightarrow$ High confidence output.
    \item \textbf{Inconsistent}: Paths diverge significantly $\rightarrow$ Trigger RAG explanation of discrepancy (e.g., ``LLM's cross-disciplinary intuition vs. DEA's local analysis'').
\end{itemize}

For submission recommendation, WuYa outputs a ranked list of candidate journals with:
\begin{itemize}
    \item Match score based on tier proximity.
    \item Per-journal explanation referencing evaluation dimensions.
    \item RAG-retrieved canonical text justifying the recommendation.
\end{itemize}

For journal review, WuYa outputs:
\begin{itemize}
    \item Five-dimensional evaluation report.
    \item Estimated tier vs. journal's actual tier.
    \item Specific improvement suggestions grounded in theory.
\end{itemize}
